\section{Vigésima iteración de la Fase 5: Tamaño máximo del cartucho}
\label{sec:F5I20}

En esta iteración cerramos la ampliación MA02: AVISAR cuando el cartucho no
cabe en la memoria de la máquina externa destino (el ESP32). No es una
verificación bloqueante: el diseño permite cartuchos grandes, lo que hace
falta saber es si el hardware destino puede con él, con un aviso claro al
generar y al cargar.

\subsection{Análisis del estado actual}
\label{sec:F5I20Analisis}

\begin{itemize}
	\item El GWC ya tiene TOPES PROPORCIONALES bloqueantes (F5-I6: Sprites,
	Segmentos y SegmentosPantalla, tope 255), que protegen el FORMATO; pero
	nada compara el tamaño RESULTANTE del cartucho con la memoria del
	hardware que lo va a ejecutar.
	\item El tamaño del GWC es un dato que el sistema YA conoce en los dos
	momentos que importan: al generar (byte a byte en
	\texttt{BaseCartucho.generarGWC}, presentado en el editor por
	\texttt{informarTamanos}) y al cargar el GWC
	(\texttt{tamGWC = fichero.size} en \texttt{\_procesarGWC}).
	\item El GWS declara la máquina (estilo, resolución, botones, teclas,
	\ldots) pero no dice NADA del hardware externo que la ejecutará: la
	memoria disponible es una propiedad de CADA máquina externa, así que su
	sitio natural es el setup: dos GWS de la misma máquina pueden
	corresponder a hardware con memoria distinta.
\end{itemize}

\subsection{Diseño: el campo \texttt{memoria} del GWS}
\label{sec:F5I20Diseno}

\begin{itemize}
	\item \textbf{Campo opcional del GWS: \texttt{memoria}, en KB (entero
	positivo).} Cada máquina externa declara en su setup cuánta memoria
	tiene el ESP32 que la ejecuta. Se escribe en el JSON del GWS como
	\texttt{"memoria": 110}.
	\item \textbf{Defecto 110 KB, SIN error.} Si el GWS no declara el campo
	(o lo declara vacío), se usa 110 KB ($110 \times 1024$ bytes) sin
	mensaje de error: es el valor razonable de un ESP32 con el reparto de
	memoria que necesita el motor (el valor se documentará en el anexo de
	formatos junto al resto de campos del GWS, PC5).
	\item \textbf{Validación con cuidado:} si el valor DECLARADO no es un
	entero positivo, aviso (F5-I16) y el sistema asume el defecto; no
	bloquea, consistente con el carácter de defecto (sin error) del campo.
	\item \textbf{La ficha del Setup} incorpora el campo «Memoria (KB)»
	(placeholder 110): vacío vale 110; a la hora de serializar, solo se
	escribe en el GWS si el autor puso un valor distinto del defecto.
	\item \textbf{¿Dónde se compara? En los DOS momentos en que el GWC
	existe como tal:} (1) al GENERAR, «Generar GWS+GWB+GWI+GWC» conoce los
	bytes del GWC recién creado y compara; (2) al CARGAR un
	\texttt{.GWC}, el tamaño ya se toma del fichero (F5-I6). El aviso es
	NO bloqueante: el cartucho se genera o carga igualmente, con la pista
	de cuánto excede (p.~ej. «el cartucho (X de 110 KB) supera la memoria
	de la máquina (F5-I20)»).
	\item \textbf{Alcance:} la memoria solo afecta al GWC: el GWB y el GWI
	viven en el navegador del autor y no van al ESP32.
\end{itemize}

\subsection{Plan de cambios}
\label{sec:F5I20Plan}

\begin{itemize}
	\item \texttt{BaseSetup.js}: campo nuevo de la ficha de setup
	(«Memoria (KB)», placeholder 110); lo toma \texttt{estado()}; solo se
	serializa en el GWS si se declaró distinto del defecto; la repoblación
	muestra el campo del GWS cargado (o vacío, que vale 110).
	\item \texttt{BaseCartucho.js}: \texttt{validarSetup} acepta el campo
	si es entero positivo (o ausente); \texttt{revisarValores(F5-I16)}
	emite aviso si el valor declarado no es válido y lo devuelve para que
	el defecto 110 se aplique sin error. Constante de módulo
	\texttt{MEMORIA\_POR\_DEFECTO = 110}.
	\item \texttt{BaseDirector.js}: función pequeña reutilizable
	(\texttt{avisarTopeMemoria(bytes)}) que compara los bytes del GWC
	contra la memoria del setup (con el defecto cuando ausente); se llama
	desde \texttt{\_procesarGWC} (tras tomar \texttt{tamGWC}) y se
	exporta para quien genera.
	\item \texttt{BaseDescomponer.js}: en \texttt{generarGWIyGWC}, tras
	crear el GWC, el mismo aviso si excede (junto a los avisos del par de
	F5-I17 M2).
	\item Cabeceras de módulos al día; A5-anexo, una línea de \texttt{
	\textbackslash input} nueva.
\end{itemize}

\subsection{Pruebas}
\label{sec:F5I20Pruebas}

\begin{itemize}
	\item Caja nueva del grupo \texttt{director} (caja D10, destructiva, al
	final de la suite S2) que se registra y se REFERENCIA EN F5-I18 (la
	norma vigente): carga un GWS cuyo campo declarado es \texttt{memoria: 1}
	y tras él el \texttt{cartucho.GWC} de \texttt{pruebas/cartuchos} (su
	tamaño real supera 1 KB): la carga sale ADELANTE (no bloquea) y
	\texttt{\#estado} canta el aviso «supera la memoria»; después, un GWS
	SIN el campo y el mismo GWC cargan SIN aviso (el defecto 110 basta).
	\item Caja del grupo \texttt{gwi\_gwc} (solo consola, C8, al estilo de
	C7): generar con un setup \texttt{memoria: 1} avisa del exceso; con el
	defecto, no.
	\item S2 y S5 completas en verde; los avisos nuevos no tocan el formato
	GWC ni el orden del tick.
\end{itemize}

\subsection{Criterios de aceptación}
\label{sec:F5I20Criterios}

\begin{enumerate}
	\item Un GWS sin \texttt{memoria} usa 110 KB sin error ni aviso.
	\item Generar (o cargar) un GWC que quepa: sin aviso; uno que exceda:
	aviso claro con los dos tamaños, NO bloqueante.
	\item La ficha de Setup muestra y guarda el campo \texttt{memoria}
	cuando se declara, con el placeholder 110.
	\item S2 y S5 en verde al cerrar la iteración.
\end{enumerate}
